\cleardoublepage

\newrefsection

\chapter{外文翻译}

\sectionnonum[stayhere]{摘要}

在多模态大型语言模型（MLLM）中，视觉投影器是视觉编码器与大型语言模型（LLM）之间不可或缺的桥梁。通常，MLLM采用简单的多层感知机（MLP）通过一对一转换来保留所有视觉上下文。然而，视觉令牌存在冗余，且在处理高分辨率图像时其数量会大幅增加，这会显著降低MLLM的效率。近期部分研究引入了重采样器或抽象器来减少生成的视觉令牌数量。遗憾的是，这些方法无法捕捉更精细的细节，从而削弱了MLLM的视觉推理能力。在本研究中，我们提出了一种新型视觉投影器，该投影器采用由粗到细的方案，通过注入丰富的特征来生成精简的视觉令牌。具体而言，我们首先将视觉特征插值为低分辨率点查询，以此作为整体视觉表征的基础。随后，引入区域到点注入模块，利用高分辨率、多级别的基于区域的线索作为精细的参考键值对，使其能够被完全融入相应的局部上下文区域中。这一步骤有效地更新了粗略的点查询，将其转化为用于后续大语言模型推理的丰富版查询。大量实验表明，我们的方法将视觉令牌压缩了75\%$\sim$89\%，同时在多种基准测试中实现了相当甚至更优的性能，且效率显著更高。源代码可在 \url{https://github.com/CircleRadon/TokenPacker} 获取。

\section{引言}

随着大型语言模型（LLM）[63, 49, 3, 50, 51, 1, 22]的快速演进，多模态大型语言模型 （MLLM）[35, 33, 34, 4, 64, 9, 15, 45, 62] 在视觉-语言理解、推理和交互能力方面取得了显著提升。这一成果是通过将视觉编码器的嵌入向量投射到LLM中来实现的，从而使LLM能够感知视觉世界，其中视觉投影器在连接视觉模型与语言模型方面发挥着关键作用。

在多模态大型语言模型的框架下，大型语言模型（LLM）通常承担了绝大多数的计算成本，尤其是考虑到视觉编码器相较于大型语言模型通常规模要小得多。例如，广泛使用的CLIP-ViT-Large [44]拥有3亿个参数，这与拥有7/8亿或13亿参数的LLaMA [50]或Vicuna [52]等LLM形成了鲜明对比。因此，MLLM的效率在很大程度上取决于视觉投影器生成的视觉令牌数量。此外，视觉投影器通过将视觉特征转换为文本嵌入空间中的视觉令牌（语言模型可解读的形式），从而连接视觉模型与语言模型。因此，这些视觉令牌的质量直接影响 MLLM 的整体效能。在本研究中，我们旨在探索一种有效的 MLLM 视觉投影器，既能以高质量连接视觉编码器与 LLM，又能尽可能减少令牌数量。

当前大多数研究采用线性投影器[35, 33]或重采样器[4, 13, 57]。就线性投影器而言，MLP投影[33]通过一对一变换保留了所有视觉上下文，这虽然保留了详细信息，但也产生了冗余的令牌[6, 46]。更重要的是，在处理高分辨率图像或视频时，视觉令牌的数量会显著增加。另一研究方向中，重采样器[4]或Q-Former[27]利用一组可学习的查询显式控制视觉令牌数量，并采用交叉注意力层来强制从视觉特征中提取最相关的视觉线索。一些近期研究，例如 Abstractor [7] 或 LDP [11, 12]，利用卷积层促进视觉特征的局部交互，从而生成压缩令牌。然而，这些方法不可避免地会丢失更精细的细节信息，并牺牲了MLLMs的视觉推理能力。此外，某些方法通过简单的像素打乱[9]或邻近拼接操作[15]将视觉特征直接从序列维度转换为通道维度，以缩短序列长度。尽管保留了所有信息，但这可能会破坏视觉特征本身的结构特征。

在本研究中，我们提出了一种名为TokenPacker的创新视觉投影器，它能将精细的细节信息有效地压缩为紧凑的视觉令牌表示。我们的TokenPacker采用“粗到细”的设计，通过将丰富的高分辨率特征注入到粗略的低分辨率特征中，从而生成紧凑的视觉令牌。具体而言，我们首先将视觉编码器输出的视觉特征插值为低分辨率点查询，其中包含视觉线索的粗略且整体的特征。随后，我们引入区域到点的注入模块，该模块充分利用高分辨率、多层级的CLIP特征，提供精细的候选键值对作为参考。在此过程中，系统会引导高分辨率视觉区域的细节注入低分辨率点查询，并在局部上下文区域内进行更新。这有效地增强了粗略查询，将其转化为更丰富的查询，以供后续的大型语言模型（LLM）使用。作为扩展，我们进一步提出了一种有效的动态图像切片方案，配合我们的TokenPacker实现高效的高分辨率图像理解。

我们在多种多模态基准数据集上进行了广泛实验，以验证本方法的有效性。值得注意的是，在LLaVA-1.5 [33]上，我们的TokenPacker能够将视觉令牌数量有效减少75\%（576 vs. 144）至89\%（576 vs. 64），同时在显著提升效率的同时，仍能保持相当甚至更优的性能。如图1所示，与其他方法相比，我们的方法在准确率和效率方面均表现出更显著的优势。此外，在多种多模态任务中，我们的方法始终能提供具有竞争力的高分辨率理解性能。

\section{相关工作}

\subsection{多模态大型语言模型}

大型语言模型（LLMs）[50, 51, 63, 3, 49, 1, 22, 1] 因其在大问答和文本生成等各类语言任务中表现出的卓越能力而备受关注。这一热潮为近期多模态大型语言模型（MLLM）[62]的开发铺平了道路，该类模型将LLM与视觉编码器相结合，从而能够对多模态内容进行更丰富的理解和解读。CLIP [44]等创新模型显著缩小了语言处理与视觉任务之间的差距，推动了跨模态应用的发展。早期研究如 Flamingo [2] 和 BLIP-2 [27]，利用了海量的图文配对数据集来优化跨模态对齐，从而大幅提升了学习效率。这一改进标志着 MLLM 领域的重要进展，通过同时处理文本和图像，拓展了应用范围。近年来，多种 MLLM 模型崭露头角。值得关注的开源模型包括LLaVA系列[35, 33, 34]、MiniGPT-4[64]、Qwen-VL[4]、CogVLM[53]、Shikra[8]、InternLM-XComposer[14]等[10, 39]。专有商业 MLLM 的出现标志着该领域格局的重大转变，例如 OpenAI 的 GPT-4V [43] 和 Google 的 Gemini 系列 [48, 45]。这些进展凸显了该领域 MLLM 多样且不断扩展的格局，这对 AGI 的发展格局产生了显著影响。

\subsection{MLLM中的可视化投影仪}

视觉投影器在连接视觉模型与语言模型方面发挥着基础性作用，它将视觉编码器输出的视觉信号映射到大型语言模型（LLM）的空间中。当前的方法主要可分为两类。一类是通过多层感知机（MLP）进行的线性投影[35, 33]。MLP投影可通过一对一变换保留所有视觉上下文，从而通过冗余令牌保留详细信息[6, 46]。该方法面临的一个关键问题是视觉令牌数量的显著增加，特别是在处理高分辨率图像或视频时。为解决这一问题，另一条研究路线专注于减少视觉令牌数量以提升MLLM的效率。Resampler [4] 或 Q-Former [27] 采用可学习的查询来显式控制视觉令牌数量，并通过交叉注意力层迫使从视觉特征中提取最相关的视觉线索。基于 Resampler，Yu 等[59]提出了查询提案网络（QPN），用于生成初始查询并执行多级交叉注意力。一些近期工作，例如 Abstractor [7] 和 LDP [11, 12]，采用卷积层来促进视觉特征的局部交互并生成压缩令牌。然而，这些方法不可避免地会遗漏精细的细节信息，从而削弱了MLLM的视觉推理能力。此外，部分研究通过简单的像素打乱[9]或邻近拼接[15]操作，将视觉特征直接从长度维度转换为通道维度以减少视觉令牌数量。尽管所有信息均得以保留，但这可能破坏视觉特征本身的内在特性。近期研究[41]对常用投影器进行了实证研究，结论是其类型影响微乎其微。与这些发现相反，本文提出了一种名为TokenPacker的新颖且有效的视觉投影器。

\subsection{利用大型语言模型实现高精度理解}

大多数大规模语言模型（MLLMs）通常采用 CLIP-ViT [44] 作为视觉编码器来捕捉视觉信息。然而，视觉编码器受限于低分辨率输入（如 224×224 或 336×336），这阻碍了 MLLMs 有效处理需要更精细细节的任务，例如密集 OCR、人群计数以及小物体的视觉定位。为克服这一限制，部分方法[54, 19, 60, 28, 39]直接采用了能高效支持高分辨率输入的视觉编码器（如SAM编码器[24]或ConvNeXt[38]），以捕捉更精细的视觉线索。与这些方法不同，本文引入了图像块裁剪策略，将高分辨率图像分割为多个图像块。随后对这些图像片段分别进行处理，从而获得整张高分辨率图像的视觉嵌入。部分研究[32, 30]首先将输入图像缩放至可处理尺寸，并采用滑动窗口将图像分割为统一尺寸的片段（例如224×224）。虽然这些方法将原始分辨率转换为固定的正方形尺寸，但这可能会导致视觉内容出现模糊或失真。为缓解这一问题，多项研究[56, 34, 16, 9]在调整图像尺寸时，采用与输入图像相似的宽高比，而非拘泥于固定的正方形比例。

\section{方法}

在本节中，我们将首先回顾标准MLLM的整体框架，该模型针对给定的多模态输入生成遵循指令的响应（第3.1节）。随后，我们将介绍名为TokenPacker的高效视觉投影器，该组件专为连接视觉编码器与LLM而设计，旨在生成紧凑的视觉令牌表示，以供后续的LLM处理使用（第3.2节）。最后，我们提出了一种动态图像切片方案，该方案支持任意宽高比的输入图像，且填充内容极少。通过集成TokenPacker，我们的方法能够以显著的计算效率实现精细的高分辨率图像理解（第3.3节）。

\subsection{重新审视多模态大型语言模型（MLLM）}

多模态大型语言模型（MLLMs）的目标是开发一种复杂的模型，使其能够根据多模态输入（包括视觉和文本数据）生成符合给定指令的响应。MLLMs通常由三个关键组件构成：1) 视觉编码器FI:它将输入图像Iimg$\in$RH×W×3转换为一组具有独特特征的视觉嵌入向量Iv$\in$RN×C.它始终采用广泛使用的 CLIP-ViT-L/14 作为其骨干网络，补丁大小 P 为 14，其中 N = HW/P 2 表示视觉嵌入的数量。2) 视觉投影器ΓI→T: 该组件将视觉嵌入Iv转换为文本嵌入空间 T 中的视觉令牌 Tv，其维度适合后续的语言模型。3) 大型语言模型 LLMΦ(Tv,Tt):它同时接收视觉令牌 Tv 和文本令牌 Tt，并通过自回归方式生成连贯的响应。对于长度为 L 的响应序列，生成符合上下文目标的答案Y={yi}Li=1的概率可通过以下公式计算：LLp(Y|Tv,Tt)=Yp(yi|Tv,Tt,<i,Y<i).(1)

在此典型的MLLM框架中，计算和内存需求主要取决于拥有大量参数的LLM Φ(Tv,Tt)。需要强调的是，LLM Φ(Tv,Tt)的计算开销通常会随着其输入令牌数量的增加而呈二次增长。这凸显了输入令牌数量对框架整体效率的显著影响。视觉投影器将 N 个视觉嵌入 Iv 转换为 M 个视觉令牌 Tv。因此，减少视觉令牌的数量是提升 LLM 效率的关键方法，即 M < N。

\subsection{TokenPacker：一款高效的可视化投影工具}

视觉投影器在将 N 个视觉嵌入 Iv 转换为 M 个视觉令牌 Tv 并输入到 LLM 之前起着至关重要的作用。如图 2 所示，我们引入了一个高效的视觉投影器——TokenPacker，它使用尽可能少的令牌将视觉编码器与语言模型连接起来。我们的TokenPacker架构采用由粗到细的框架。具体而言，我们在基于CLIP的视觉编码器最后一个Transformer层之前，通过双线性插值（缩放因子s）对视觉特征Iv $\in$ RN×C进行下采样，将其转换为低分辨率视觉嵌入I′ $\in$ RM×C，其中M = N/s²。因此，视觉令牌数 Mv′ 可通过下采样率 s 进行控制。低分辨率 Iv 可视为原始高分辨率视觉特征的粗略表示，其中低分辨率 I′ 的每个像素对应高分辨率 Iv 中的特定 (s×s) 子区域。随后，我们构建点-区域对，即 I′ $\in$ R1×M×C 中的每个像素与 Iv $\in$ Rs ×M×C 中的子区域，并旨在将高分辨率子区域的详细信息注入到具有粗略表示的每个像素中。为了完成这一过程，我们设计了一个注入模块，该模块有效地执行区域到点的信息注入，以增强和更新低分辨率表示。

具体而言，我们将低分辨率图像 I′ $\in$ R1×M×C 作为点级查询，并将 Iv $\in$ Rs ×M×C 作为区域级候选键值进行参照。区域到点的信息注入通过多层感知器（MLP）层后的点到区域交叉注意力操作实现，使低分辨率查询充分吸收细粒度的键值对，并更新为紧凑且增强的视觉令牌 Tv。此外，我们利用多级视觉特征作为更丰富的参考键值对。正如先前的研究 [23] 所示，CLIP 编码器的不同层对不同模式表现出不同的偏好。浅层特征包含详细的低级信息，而深层特征则在语义理解方面更具优势。多层区域到点注入过程有助于将多层丰富的高分辨率信息注入低分辨率查询中，使其足以充当视觉令牌。因此，我们的方法能够生成更优质的视觉令牌，同时将视觉令牌的总数减少到视觉嵌入的 1/√2。

\section{实验}

在本节中，我们将首先介绍实验设置的详细情况。随后，我们将针对多种多模态测试平台，将我们的方法与主流方法进行性能对比。本节最后，将呈现消融分析及定性结果。

\subsection{实施细节}

在本研究中，我们基于 LLaVA-1.5 [33] 实现了我们的方法。具体而言，我们采用 CLIP-ViT-L/14-336px [44] 作为视觉编码器（默认分辨率为 336×336），并选用 Vicuna-7/13B 模型 [63] 作为大型语言模型（LLM）。我们采用两阶段训练范式，包括预训练阶段和指令微调阶段。为确保训练效率，我们在两个阶段中均保持视觉编码器固定，同时专注于优化我们提出的TokenPacker。

与此同时，LLM的优化仅在指令微调阶段进行。我们通过调整TokenPacker中的下采样率s $\in$ {2, 3, 4}来控制生成的视觉令牌数量。在针对高分辨率图像的动态切片方案中，我们将Ng设为9或16用于模型训练和评估，以支持1344×1344、5376×336等多种分辨率。在式(4)中，默认将 β 设为 0.1。与 [33] 一致，我们利用 AdamW 优化器配合余弦学习率曲线，对所有模型进行一轮训练。预训练阶段和指令调优阶段的初始学习率分别设为 1e−3 和 2e−5。模型在 8 块 NVIDIA A100 GPU 上进行训练。

\subsection{数据集与基准测试}

为了进行公平的比较，我们首先在 CC-558K 数据集 [35] 上进行实验，以训练我们的 TokenPacker，从而在第一阶段完成模态对齐。第二阶段则参照 LLaVA-1.5 [33] 的做法，使用 665K 混合数据集 [35] 进行指令微调。为了取得具有竞争力的性能，我们随后采用了Mini-Gemini [28]中整理的更多高质量训练样本，第一阶段约120万，第二阶段约150万。此外，我们在一系列广泛使用的基准测试上进行了全面评估，以评估我们提出的模型的多模态理解和推理能力。本研究采用的基准测试包括：1) 通用视觉问答基准，如 VQAv2 [17]、GQA [20]、VizWiz [18]； 2) 与OCR相关的基准，如VQAT（TextVQA）[47]、OCRBench（OCRB）[37]和DocumentVQA（DocVQA）[40]； 3) 幻觉评估基准，如POPE [29]；4) 综合性基准，如MMB (MMBench) [36]、MM-Vet [58]和MMMU [61]。

\subsection{主要结果}

标准分辨率。我们首先利用LLaVA-1.5 [33]中的数据，在标准分辨率设置下考察所提出的TokenPacker的有效性。我们将本方法与当前领先的方法进行比较，包括MobileVLM V2 [12]、Shikra [8]、IDEFICS [21]、Qwen-VL [4]、InstructBLIP [13]，以及视觉令牌数量较少的LLaVA-PruMerge [46]。我们采用了六个流行的基准测试，包括全面的MMBench和MM-Vet，以及通用的VQA相关任务VizWiz、VQAv2、GQA，以及用于评估幻觉生成能力的POPE，以进行全面的性能评估。如表1所示，我们的方法在MMBench、VizWiz和POPE基准测试中分别展现了优异的性能。与基线模型 LLaVA-1.5 相比，我们提出的 TokenPacker 作为视觉投影器，将视觉令牌数量减少了 75\%（从 576 降至 144），同时在 MMBench 上将性能指标提升了 +0.8\%/+0.3\%， 在VizWiz上两项指标均提升2.0\%，并在POPE上分别提升1.1\%和1.5\%（分别采用Vicuna-7B和Vicuna-13B大语言模型）。尽管在VQAv2和GQA等图像问答基准测试中性能略有下降，但我们的方法仍全面提升了相对于LLaVA-1.5 [33]的平均性能， 使用 Vicuna-7B 和 Vicuna-13B 模型时，分别提升 +0.8\% 和 +0.1\%，且每秒处理量（TPS）约为前者的 5 倍（4.9 对比 24.9，详见表 3）。此外，在绝大多数基准测试中，我们的方法都超越了Qwen-VL-Chat [4]、InstructBLIP [13]和MobileVLM V2 [12]等先前方法，尽管这些方法拥有更丰富的训练数据。

此外，我们还对比了本方法与使用较少视觉令牌的先前领先方法。具体而言，我们将令牌数分别设置为 64（占 576 的 11\%）和 36（占 576 的 6\%）。可以看出，本方法在三个基准测试中均以较大优势超越了这些方法。例如，我们观察到，在MMBench上比近期采用Vicuna-13B的LLaVA-PruMerge [46]高出3.9\%，在VQAv2上高出3.5\%。这些结果证实了我们的TokenPacker的有效性，突显了其对增强视觉令牌表示及整体性能的积极影响。

高分辨率。我们将包括TokenPacker和动态图像切片方案在内的方法应用于LLaVA-1.5（命名为LLaVA-TokenPacker-HD），以实现高分辨率图像理解。模型训练采用Mini-Gemini [28]中组织的270万张数据。我们将Ng设为9和16，以支持最大输入分辨率分别为1088×1088和1344×1344。下采样比s设为2、3或4，以控制从每个图像补丁中提取的视觉令牌数量。我们将本方法与现有高分辨率MLLM方法（如OtterHD [26]、SPHINX-2k [32]、Monkey [30]、面向文档的UReader [56]和TextHawk [59]）以及较新的LLaVA-UHD [55]、LLaVA-NeXT [34]、 Mini-Gemini-HD [28]。表2报告了在九个流行基准上的对比结果，包括与OCR相关的VQAT、OCRB和DocVQA，以及综合性的MMB、MMMU和MME，还有通用的VQA相关基准VQAv2、VizWiz和POPE。可以看出，当输入分辨率设置为 1344×1344（约 1393 个视觉令牌）时，采用 Vicuna-13B 作为 LLM 的我们的方法在 VQAT 上实现了 70.6\% 的 OCR 相关最先进性能，在 OCRBench 上实现了 521 的成绩。这些令人鼓舞的结果可归因于高分辨率图像为MLLM提供了更多的视觉令牌，从而能够精确识别复杂且精细的光学字符或物体。然而，在MMMU和MME等综合性基准测试中，我们的方法在1088×1088的较低分辨率下表现最佳。此外，即使视觉令牌数量仅约619个，我们的方法在MMMU、MME和VizWiz上的得分仍分别达到38.2\%、1577/353和61.0\%，位列第二。这些结果表明，即使令牌数量减少，MLLM在综合性基准测试和VQA相关任务中仍能展现出稳健的性能。这些结果证明了在多样化的多模态任务中利用原生高分辨率图像具有关键性作用，并突显了我们提出的TokenPacker的有效性。图4展示了在代表性场景中的定性比较。

\subsection{消融结果}

我们进一步深入探讨了详细的消融实验，以分析我们方法中每个组件的有效性。所有消融实验均采用Vicuna-7B作为大语言模型（LLM），并使用LLaVA-1.5 [33]中的数据进行模型训练。

多种视觉投影器。首先，我们在LLaVA-1.5模型基础上，将我们提出的TokenPacker与多种先前的视觉投影器进行比较，包括直接平均池化、Resampler [4]、C-Abstractor [7]、Pixel-Shuffle [9]以及最新的LDP-v2 [12]。对于平均池化方法，我们直接通过平均池化对视觉编码器输出的特征图进行插值，随后使用多层感知机（MLP）生成视觉令牌。为确保公平比较，我们将 LLaVA-1.5 中的原始 MLP 替换为各种投影器，并保持相同设置。为反映推理速度，我们采用每秒令牌数（TPS）指标来评估 LLM 的吞吐量。表 3 报告了对比结果。与 MLP 投影器相比，所有其他视觉投影器均有效减少了视觉令牌数量，同时显著提升了推理速度（约 5 TPS 对比 25 TPS）。平均池化在参数更少的情况下实现了最佳的推理速度。我们的视觉投影器在 144 个令牌下，相较于 576 个令牌的 MLP，平均性能提升了 0.8\%。特别是在 MM-Vet [58]、POPE [29] 和 VizWiz [18] 基准测试中，我们的方法分别比基于 MLP 的方法高出 1.9\%、1.1\% 和 2.0\%。与其他方法相比，我们的方法比此前最佳方法 LDP-v2 [12] 高出 1.6\%。在 64 个令牌的场景下，我们的方法实现了 61.9\% 的平均性能，与基于 MLP 的方法（61.9\% 对比 62.0\%）持平，并比 LDP-v2 [12] 高出 1.3\%。这些结果证明了 TokenPacker 相较于先前方法的有效性。

\section{结论与局限性}

在本研究中，我们为MLLM提出了一种新型视觉投影器——TokenPacker。我们的方法采用由粗到细的设计，能够将丰富的高分辨率图像特征有效地浓缩为紧凑的视觉令牌。作为扩展，我们进一步提出了一种有效的动态图像划分方案，以实现高效的高分辨率图像理解。我们已在多种基准测试上进行了广泛的实验，以验证我们方法的有效性。值得注意的是，我们的TokenPacker在LLaVA-1.5数据集上可将视觉令牌数量有效减少75\%$\sim$89\%，同时在保持相当甚至更优性能的同时，实现了显著更高的效率。

局限性。我们的TokenPacker通过将视觉令牌压缩高达89\%展现了出色的性能，但并非完全无损。具体而言，当令牌数量减少至32个（6\%）或更少时，性能会出现明显下降。我们将致力于推进研究，开发出仅需极少量令牌的更精密视觉投影器，以实现基于MLLM的高效视觉理解。
